\chapter{Probabilités et variables aléatoires}\label{ch:g11:prob}

La \omterm{def:g10:proba:distribution}{probabilité} modélise des expériences dont le résultat est incertain~:
un lancer de dé, une carte tirée, un jeu joué. Ce chapitre met en place
le vocabulaire sur des ensembles finis, puis introduit le personnage
central de toute la théorie des \omterm{def:g10:proba:distribution}{probabilités}~: la \emph{\omterm{def:g11:prob:rv}{variable
aléatoire}} et son \emph{\omterm{def:g11:prob:expectation}{espérance}}. Ces notions sont reprises et
approfondies dans le \cref{ch:g12:randvar}, et les \omterm{def:g10:proba:distribution}{probabilités}
conditionnelles suivent dans le \cref{ch:g12:condprob}.

\section{Probabilité sur un ensemble fini}

\begin{definition}[Modèle probabiliste]\label{def:g11:prob:model}
Une expérience à un nombre fini d'issues est modélisée par son
\emph{univers}\index{univers} $\Omega = \{\omega_1, \dots, \omega_n\}$
(l'ensemble des issues) muni d'une \emph{\omterm{def:g10:proba:distribution}{probabilité}} $\P$ attribuant à
chaque issue $\omega_i$ un nombre $p_i \geq 0$, avec
$p_1 + \dots + p_n = 1$. Un \emph{événement}\index{événement} est une
partie $A \subseteq \Omega$, et
\[
\P(A) = \sum_{\omega_i \in A} p_i .
\]
Lorsque toutes les issues sont équiprobables ($p_i = \frac1n$~: le
\emph{modèle uniforme}\index{modèle uniforme}),
\[
\P(A) = \frac{\text{nombre d'issues dans } A}
{\text{nombre d'issues dans } \Omega}.
\]
\end{definition}

\begin{example}\label{ex:g11:prob:twodice}
Lancer deux dés distincts~: $\Omega$ est l'ensemble des $36$ couples
ordonnés $(i, j)$ avec $1 \leq i, j \leq 6$, tous équiprobables.
L'\omterm{def:g11:prob:model}{événement} $S$ = «~la somme est $7$~» contient les $6$ couples
$(1,6), (2,5), (3,4), (4,3), (5,2), (6,1)$, donc
$\P(S) = \frac{6}{36} = \frac16$. L'\omterm{def:g11:prob:model}{événement} «~la somme est $12$~»
contient seulement $(6,6)$~: \omterm{def:g10:proba:distribution}{probabilité} $\frac{1}{36}$. Choisir le
\emph{bon} \omterm{def:g11:prob:model}{univers} --- couples ordonnés, pas non ordonnés --- est ce qui
rend le \omterm{def:g11:prob:model}{modèle uniforme} applicable.
\end{example}

\begin{proposition}[Règles de la probabilité]\label{prop:g11:prob:rules}
Pour des \omterm{def:g11:prob:model}{événements} $A, B \subseteq \Omega$~:
\[
\P(\emptyset) = 0, \qquad \P(\Omega) = 1, \qquad
\P(\bar A) = 1 - \P(A),
\]
\[
\P(A \cup B) = \P(A) + \P(B) - \P(A \cap B),
\]
où $\bar A$ est le \emph{\omterm{def:g10:proba:operations}{complémentaire}} de $A$ (toutes les issues hors
de $A$). En particulier, $\P(A \cup B) = \P(A) + \P(B)$ lorsque $A$ et
$B$ sont \emph{\omterm{def:g10:proba:operations}{incompatibles}} ($A \cap B = \emptyset$).
\end{proposition}

\begin{proof}
L'\omterm{def:g11:prob:model}{événement} vide ne contient aucune issue (somme vide), et $\Omega$ les
contient toutes (somme totale $1$). Chaque issue est dans exactement
l'un de $A$, $\bar A$, donc $\P(A) + \P(\bar A) = 1$. Pour la \omterm{def:g10:numbers:interunion}{réunion}~:
sommer les $p_i$ sur $A$ puis sur $B$ compte deux fois les issues de
$A \cap B$~; soustraire $\P(A \cap B)$ corrige le double comptage.
\end{proof}

\begin{example}
Dans une classe, $60\,\%$ étudient la musique ($M$), $45\,\%$ font un
sport ($S$), et $25\,\%$ font les deux. La proportion qui fait au moins
l'un des deux est $\P(M \cup S) = 0.60 + 0.45 - 0.25 = 0.80$, donc
$20\,\%$ ne font ni l'un ni l'autre (\omterm{def:g10:proba:operations}{complémentaire}).
\end{example}

\section{Variables aléatoires}

\begin{definition}[Variable aléatoire, loi]\label{def:g11:prob:rv}
Une \emph{variable aléatoire}\index{variable aléatoire} sur $\Omega$ est
une \omterm{def:g11:func:function}{fonction} $X$ attribuant un nombre à chaque issue. Sa
\emph{loi}\index{loi} est le tableau de ses valeurs possibles
$x_1, \dots, x_k$ avec leurs \omterm{def:g10:proba:distribution}{probabilités}
\[
\P(X = x_i) = \P\bigl(\{\omega : X(\omega) = x_i\}\bigr),
\qquad
\sum_{i=1}^k \P(X = x_i) = 1 .
\]
\end{definition}

\begin{example}\label{ex:g11:prob:game}
Un jeu~: payer $2$ euros, lancer un dé équilibré, recevoir la valeur
affichée si elle est au moins $5$, rien sinon. Soit $G$ le
\emph{gain} (reçu moins payé). Les issues $1, 2, 3, 4$ donnent
$G = -2$~; l'issue $5$ donne $G = 3$~; l'issue $6$ donne $G = 4$. La \omterm{def:g11:prob:rv}{loi}
de $G$~:
\begin{center}
\begin{tabular}{c|ccc}
$g$ & $-2$ & $3$ & $4$\\
\hline\rule{0pt}{11pt}%
$\P(G = g)$ & $\dfrac46$ & $\dfrac16$ & $\dfrac16$
\end{tabular}
\end{center}
\end{example}

\begin{method}[Construire un tableau de loi]\label{met:g11:prob:table}
Lister les valeurs possibles de $X$~; pour chaque valeur, rassembler les
issues qui la produisent et additionner leurs \omterm{def:g10:proba:distribution}{probabilités}~; vérifier
que la ligne des \omterm{def:g10:proba:distribution}{probabilités} somme à $1$.
\end{method}

\section{Espérance}

\begin{definition}[Espérance]\label{def:g11:prob:expectation}
L'\emph{espérance}\index{espérance} (ou \omterm{def:g11:stat:mean}{moyenne}) de $X$ est la \omterm{def:g11:stat:mean}{moyenne} de
ses valeurs pondérées par leurs \omterm{def:g10:proba:distribution}{probabilités}~:
\[
\E(X) = \sum_{i=1}^{k} \P(X = x_i)\; x_i .
\]
\end{definition}

\begin{remark}
$\E(X)$ est ce que serait la \omterm{def:g11:stat:mean}{moyenne} d'un grand nombre de répétitions
indépendantes de l'expérience --- la version précise de cette affirmation
est la loi des grands nombres, prouvée dans le \cref{ch:g12:sums}. Noter
l'analogie formelle avec la \omterm{def:g11:stat:mean}{moyenne} d'un jeu de données avec effectifs
(\cref{def:g11:stat:mean})~: les \omterm{def:g10:proba:distribution}{probabilités} jouent le rôle des
\omterm{def:g10:stats:series}{fréquences} relatives $\frac{n_i}{n}$.
\end{remark}

\begin{example}\label{ex:g11:prob:gameexp}
Pour le jeu de l'\cref{ex:g11:prob:game}~:
\[
\E(G) = \frac46 \times (-2) + \frac16 \times 3 + \frac16 \times 4
= \frac{-8 + 3 + 4}{6} = -\frac16 \approx -0.17 .
\]
En \omterm{def:g11:stat:mean}{moyenne} le joueur perd environ $17$ centimes par partie~: le jeu est
défavorable. Un jeu est dit \emph{équitable} lorsque le gain espéré est
$0$.
\end{example}

\begin{proposition}[Linéarité]\label{prop:g11:prob:linearity}
Pour tous réels $a, b$~:
\[
\E(aX + b) = a\,\E(X) + b .
\]
\end{proposition}

\begin{proof}
La variable $aX + b$ prend la valeur $a x_i + b$ avec la \omterm{def:g10:proba:distribution}{probabilité}
$\P(X = x_i)$, donc
\[
\E(aX + b) = \sum_i \P(X = x_i)(a x_i + b)
= a \sum_i \P(X = x_i)\,x_i + b \underbrace{\sum_i \P(X =
x_i)}_{=\,1} = a\,\E(X) + b .
\]
\end{proof}

\section{Variance et écart type}

\begin{definition}[Variance]\label{def:g11:prob:variance}
La \emph{variance}\index{variance} de $X$ mesure la dispersion de ses
valeurs autour de l'\omterm{def:g11:prob:expectation}{espérance} $m = \E(X)$~:
\[
\V(X) = \E\bigl((X - m)^2\bigr)
= \sum_{i=1}^{k} \P(X = x_i)\,(x_i - m)^2,
\qquad
\sigma(X) = \sqrt{\V(X)} .
\]
\end{definition}

\begin{proposition}[Formule de König--Huygens]\label{prop:g11:prob:konig}
\[
\V(X) = \E(X^2) - \E(X)^2 .
\]
\end{proposition}

\begin{proof}
Développer $(x_i - m)^2 = x_i^2 - 2m\,x_i + m^2$ dans la somme~:
\[
\V(X) = \sum_i \P(X = x_i)\,x_i^2
- 2m \sum_i \P(X = x_i)\,x_i + m^2\sum_i \P(X = x_i)
= \E(X^2) - 2m^2 + m^2 ,
\]
qui est $\E(X^2) - m^2$. C'est la version \omterm{def:g11:prob:rv}{variable aléatoire} de
l'identité prouvée pour les jeux de données dans la
\cref{prop:g11:stat:shortcut}.
\end{proof}

\begin{proposition}[Transformation affine]\label{prop:g11:prob:affine}
$\V(aX + b) = a^2\,\V(X)$ et $\sigma(aX + b) = \abs a\,\sigma(X)$.
\end{proposition}

\begin{proof}
Par linéarité, $aX + b$ s'écarte de son \omterm{def:g11:prob:expectation}{espérance} $a m + b$ de
$(aX + b) - (am + b) = a(X - m)$~: chaque écart est multiplié par $a$,
donc chaque écart au carré par $a^2$, et leur \omterm{def:g11:stat:mean}{moyenne} pondérée aussi. Le
décalage $b$ a disparu --- décaler une variable ne change pas sa
dispersion.
\end{proof}

\begin{example}\label{ex:g11:prob:variance}
Pour le jeu de l'\cref{ex:g11:prob:game}~:
$\E(G^2) = \frac46 \times 4 + \frac16 \times 9 + \frac16 \times 16
= \frac{16 + 9 + 16}{6} = \frac{41}{6}$, donc
\[
\V(G) = \frac{41}{6} - \left(-\frac16\right)^2
= \frac{41}{6} - \frac{1}{36} = \frac{245}{36} \approx 6.8,
\qquad
\sigma(G) \approx 2.6 .
\]
La perte de $0.17$ par partie en \omterm{def:g11:stat:mean}{moyenne} s'accompagne de fluctuations de
taille typique $2.6$ euros --- l'\omterm{def:g11:prob:expectation}{espérance} sans l'\omterm{def:g11:stat:variance}{écart type} ne raconte
que la moitié de l'histoire.
\end{example}

\begin{omfigure}
\begin{tikzpicture}
\begin{axis}[
  width=8.6cm, height=5.2cm,
  ybar, bar width=14pt,
  xlabel={$g$}, ylabel={$\P(G = g)$},
  xmin=-3.5, xmax=5.5, ymin=0, ymax=0.8,
  xtick={-2,3,4}, ytick={1/6, 4/6},
  yticklabels={$\tfrac16$,$\tfrac46$},
  tick label style={font=\small},
  label style={font=\small}]
  \addplot[fill=omDef!40, draw=omDef] coordinates {
    (-2,0.6667) (3,0.1667) (4,0.1667)};
  \draw[omThm, dashed, thick] (axis cs:-0.1667,0) --
    (axis cs:-0.1667,0.75);
  \node[omThm, font=\small, anchor=south west] at (axis cs:-0.05,0.62)
    {$\E(G)$};
\end{axis}
\end{tikzpicture}

{\small La \omterm{def:g11:prob:rv}{loi} du gain $G$ de l'\cref{ex:g11:prob:game}, et son \omterm{def:g11:prob:expectation}{espérance}
$\E(G) = -\frac16$ (en pointillés)~: le point d'équilibre des masses de
\omterm{def:g10:proba:distribution}{probabilité}.}
\end{omfigure}

\section{Exercices}

\begin{exercise}[$\star$]\label{exo:g11:prob:1}
Une carte est tirée d'un jeu standard de $52$ cartes. Calculer la
\omterm{def:g10:proba:distribution}{probabilité} de tirer~: un cœur~; un roi~; un cœur ou un roi~; ni un cœur
ni un roi.
\end{exercise}

\begin{exercise}[$\star$]\label{exo:g11:prob:2}
Deux dés équilibrés sont lancés et leur somme $S$ est enregistrée. En
utilisant le modèle à $36$ issues de l'\cref{ex:g11:prob:twodice},
calculer $\P(S = 6)$, $\P(S \geq 10)$ et $\P(S \text{ est pair})$.
\end{exercise}

\begin{exercise}[$\star$]\label{exo:g11:prob:3}
Une \omterm{def:g11:prob:rv}{variable aléatoire} a pour \omterm{def:g11:prob:rv}{loi}
\begin{center}
\begin{tabular}{c|cccc}
$x$ & $-1$ & $0$ & $2$ & $5$\\
\hline
$\P(X=x)$ & $0.3$ & $0.2$ & $0.4$ & $p$
\end{tabular}
\end{center}
Trouver $p$, puis calculer $\E(X)$.
\end{exercise}

\begin{exercise}[$\star$]\label{exo:g11:prob:4}
$X$ satisfait $\E(X) = 3$ et $\V(X) = 4$. Donner l'\omterm{def:g11:prob:expectation}{espérance}, la \omterm{def:g11:prob:variance}{variance}
et l'\omterm{def:g11:stat:variance}{écart type} de $Y = 2X - 5$.
\end{exercise}

\begin{exercise}[$\star\star$]\label{exo:g11:prob:5}
Une roue de la fortune a $10$ secteurs égaux~: $5$ affichent $0$ euro,
$3$ affichent $2$ euros, $2$ affichent $5$ euros. Jouer coûte $2$ euros.
Donner la \omterm{def:g11:prob:rv}{loi} du gain $G$ (gains moins mise), calculer $\E(G)$, et
décider si le jeu est favorable.
\end{exercise}

\begin{exercise}[$\star\star$]\label{exo:g11:prob:6}
Une urne contient $3$ boules rouges et $2$ bleues~; deux boules sont
tirées \emph{sans} remise. Soit $X$ le nombre de boules rouges tirées.
Construire la \omterm{def:g11:prob:rv}{loi} de $X$ (énumérer les $20$ tirages ordonnés, ou
utiliser un \omterm{met:g10:proba:tree}{arbre}), et calculer $\E(X)$.
\end{exercise}

\begin{exercise}[$\star\star$]\label{exo:g11:prob:7}
Pour la roue de l'\cref{exo:g11:prob:5}, calculer $\V(G)$ et $\sigma(G)$.
\end{exercise}

\begin{exercise}[$\star\star$]\label{exo:g11:prob:8}
Une compagnie d'assurance vend une police pour $80$ euros. Avec
\omterm{def:g10:proba:distribution}{probabilité} $0.05$ le client réclame $1000$ euros, avec \omterm{def:g10:proba:distribution}{probabilité}
$0.01$ il réclame $5000$ euros, sinon rien. Soit $B$ le bénéfice de la
compagnie sur une police. Donner la \omterm{def:g11:prob:rv}{loi} de $B$, calculer $\E(B)$, et
interpréter.
\end{exercise}

\begin{exercise}[$\star\star$]\label{exo:g11:prob:9}
Une question à choix multiples a $4$ réponses, une seule correcte. Une
réponse correcte rapporte $+1$~; pour décourager le hasard, une réponse
fausse rapporte $x < 0$. Un élève répond au hasard. Pour quelle valeur
de $x$ le score espéré de l'élève est-il nul~?
\end{exercise}

\begin{exercise}[$\star\star$]\label{exo:g11:prob:10}
Deux pièces équilibrées sont lancées. Anne propose~: «~si les deux
pièces coïncident, je vous paie $1$ euro~; si elles diffèrent, vous me
payez $1$ euro.~» Calculer le gain espéré de chaque joueur. Le jeu
est-il équitable~? Même question avec trois pièces, où Anne paie
$2$ euros si les trois coïncident.
\end{exercise}

\begin{exercise}[$\star\star\star$]\label{exo:g11:prob:11}
Une loterie vend $1000$ billets à $5$ euros chacun. Un billet gagne
$2000$ euros, cinq billets gagnent $200$ euros, vingt billets gagnent
$25$ euros.
\begin{enumerate}
  \item Donner la \omterm{def:g11:prob:rv}{loi} du revenu total de l'organisateur moins les lots,
        et du gain $G$ d'un seul joueur.
  \item Calculer $\E(G)$ et le profit total espéré de l'organisateur.
        Vérifier que les deux réponses sont cohérentes.
\end{enumerate}
\end{exercise}

\section{Problème~: la partie interrompue}

\begin{problem}[{Devoir du week-end --- Pascal, Fermat et le partage
équitable d'un pari interrompu~: la naissance de l'espérance, la
magie de la linéarité, et ce que l'espérance ne voit pas}]
\label{pb:g11:prob:1}
En 1654, le chevalier de Méré --- le même joueur dont les paris aux
dés ouvraient le problème de \omterm{def:g10:proba:distribution}{probabilités} du volume précédent ---
posa à Blaise Pascal une question plus profonde~: deux joueurs de
force égale doivent interrompre leur série au score de $5$ à $3$,
six victoires emportant les $64$ pistoles de l'enjeu. \emph{Comment
partager équitablement l'enjeu~?} Les lettres échangées entre Pascal
et Fermat pour y répondre ont fondé la notion au cœur de ce
chapitre~: l'\omterm{def:g11:prob:expectation}{espérance} (\cref{def:g11:prob:expectation}). Vous
retrouverez leur réponse par les deux voies, puis rencontrerez le
tour de passe-passe de la linéarité et l'unique angle mort de
l'\omterm{def:g11:prob:expectation}{espérance}.

\medskip
\noindent\textbf{Partie I --- Aisance avec l'\omterm{def:g11:prob:expectation}{espérance}.}
\begin{enumerate}
  \item On lance un dé équilibré~; le gain est de (face $- 3$)
        euros. Dresser le tableau de la \omterm{def:g11:prob:rv}{loi} du gain $G$
        (\cref{met:g11:prob:table}) et calculer $\E(G)$.
  \item Calculer $V(G)$ avec la formule raccourcie
        (\cref{prop:g11:prob:konig}) puis $\sigma(G)$ (deux
        décimales).
  \item Sans dresser de nouveau tableau, donner $\E(2G - 1)$ et
        $V(2G - 1)$ (\cref{prop:g11:prob:affine}).
  \item Un stand fait payer $m$ euros pour lancer deux dés et vous
        verse la somme obtenue en euros. Quel est le prix
        \emph{équitable} $m$~? Le stand fait payer $8$~: quelle est
        votre perte espérée par partie~?
  \item Un ordinateur portable de $800$ euros a $5\,\%$ de risque
        de destruction par an. Calculer la prime d'assurance
        équitable, ainsi que le bénéfice espéré de l'assureur par
        contrat s'il fait payer $60$ euros.
\end{enumerate}

\medskip
\noindent\textbf{Partie II --- Le partage de l'enjeu.}
Score $5$--$3$, le premier à $6$ emporte les $64$ pistoles, chaque
manche étant un lancer de pièce équilibrée.
\begin{enumerate}[resume]
  \item Deux réponses fausses mais tentantes~: partager selon
        $5 : 3$ (proportionnellement aux manches gagnées), ou tout
        donner à celui qui mène. Dire en une phrase pour chacune ce
        qu'elle a d'inéquitable.
  \item La méthode de Fermat~: le joueur A a besoin d'$1$ victoire
        de plus, B de $3$~; imaginons que l'on joue exactement $3$
        manches supplémentaires quoi qu'il arrive. Énumérer les $8$
        issues équiprobables, marquer qui remporte la série dans
        chacune, et en déduire le partage équitable des $64$
        pistoles.
  \item Le point délicat que Fermat a dû défendre~: la partie
        réelle pourrait s'arrêter après une seule manche --- pourquoi
        est-il malgré tout légitime de raisonner sur les $3$ manches
        imaginaires~?
  \item La méthode de Pascal, à rebours depuis la fin~: calculer la
        part équitable de A aux scores (hypothétiques) $5$--$5$,
        puis $5$--$4$, puis $5$--$3$, en moyennant chaque fois les
        deux futurs équiprobables. Comparer avec la question 7.
  \item Même question au score $4$--$3$~: combien de manches
        imaginaires, quelles issues donnent la série à B, et quel
        est le partage équitable~?
  \item Énoncer le principe commun aux deux méthodes --- la valeur
        équitable d'un avenir incertain est son \emph{\omterm{def:g11:prob:expectation}{espérance}} ---
        et nommer deux industries modernes qui vivent entièrement de
        ce principe (la question 5 en montre une).
  \item Un contrôle de cohérence valable à tout score~: expliquer,
        à l'aide de la linéarité de l'\omterm{def:g11:prob:expectation}{espérance}
        (\cref{prop:g11:prob:linearity}), pourquoi les parts
        équitables des deux joueurs doivent toujours totaliser
        exactement les $64$ pistoles.
\end{enumerate}

\medskip
\noindent\textbf{Partie III --- Le tour de passe-passe de la
linéarité.}
\begin{enumerate}[resume]
  \item Retrouver $\E(\text{somme de deux dés}) = 7$ en une ligne
        par la linéarité (\cref{prop:g11:prob:linearity}), sans le
        tableau à $36$ cases. Pour la \emph{\omterm{def:g11:prob:variance}{variance}} de la somme,
        la règle $V(X + Y) = V(X) + V(Y)$ exige que les dés soient
        indépendants (admis)~: la calculer.
  \item La soirée aux chapeaux~: $n$ invités jettent leur chapeau
        en tas et chacun en reprend un au hasard. Soit $X$ le nombre
        d'invités qui récupèrent leur propre chapeau. Pour un invité
        fixé, quelle est la \omterm{def:g10:proba:distribution}{probabilité} de récupérer son chapeau~?
        En sommant ces $n$ chances (linéarité~!), calculer $\E(X)$.
        La réponse devrait vous surprendre~: elle ne dépend pas de
        $n$.
  \item Vérifier $\E(X) = 1$ par force brute pour $n = 2$ et
        $n = 3$ (énumérer les $2$ puis les $6$ distributions de
        chapeaux équiprobables et moyenner les comptages).
  \item Les récupérations des invités sont loin d'être
        indépendantes (si $n - 1$ invités ont retrouvé leur chapeau,
        le dernier aussi). Où exactement le calcul de la question 14
        n'a-t-il \emph{pas} eu besoin de l'indépendance~? Énoncer le
        superpouvoir de la linéarité.
\end{enumerate}

\medskip
\noindent\textbf{Partie IV --- Ce que l'\omterm{def:g11:prob:expectation}{espérance} ne voit pas.}
\begin{enumerate}[resume]
  \item Deux tickets à gratter à $5$ euros~: le ticket A fait
        gagner $10$ euros avec la \omterm{def:g10:proba:distribution}{probabilité} $\frac12$~; le ticket
        B fait gagner $1\,000$ euros avec la \omterm{def:g10:proba:distribution}{probabilité} $0.005$.
        Vérifier que les deux gains ont la même \omterm{def:g11:prob:expectation}{espérance}, puis
        calculer les deux \omterm{def:g11:prob:variance}{variances}. Quel ticket les gens
        achètent-ils réellement, et qu'achètent-ils au juste~?
  \item Le jeu de Saint-Pétersbourg~: on lance une pièce jusqu'au
        premier face~; si ce premier face survient au lancer $n$, on
        reçoit $2^n$ euros. Calculer la contribution de chaque $n$ à
        l'\omterm{def:g11:prob:expectation}{espérance}, puis l'\omterm{def:g11:prob:expectation}{espérance} elle-même. Paieriez-vous
        $1\,000$ euros pour jouer une seule fois~? Que ce heurt
        enseigne-t-il sur l'\omterm{def:g11:prob:expectation}{espérance}~?
  \item L'assurance vue du côté du client~: payer $20$ euros, ou
        affronter un risque de $1\,\%$ de perdre $1\,000$~? Comparer
        les deux \omterm{def:g11:prob:expectation}{espérances}, et expliquer pourquoi souscrire
        l'assurance peut malgré tout être rationnel (sur quoi une
        famille seule ne peut-elle pas «~faire la \omterm{def:g11:stat:mean}{moyenne}~»~?).
  \item Pour finir --- la biographie de l'\omterm{def:g11:prob:expectation}{espérance}, une phrase par
        chapitre~: née en 1654 pour partager un enjeu~; son
        superpouvoir, la linéarité, qui additionne les \omterm{def:g11:prob:expectation}{espérances}
        sans jamais demander l'indépendance (les chapeaux)~; son
        angle mort, le risque, que mesure la \omterm{def:g11:prob:variance}{variance} (les tickets,
        Saint-Pétersbourg, l'assurance)~; et son couronnement l'an
        prochain, quand la loi des grands nombres expliquera
        exactement pourquoi la banque finit toujours par gagner.
\end{enumerate}
\end{problem}
